\begin{theorem}[Binding of the GBCA specification, on a trace]
  \label{thm:aba-gbca-binding-trace}
  \lean{PLTS.ABA.GBCA.BoundTrace}\lean{PLTS.ABA.GBCA.BindingTrace}\lean{PLTS.ABA.GBCA.specInst_binding}\lean{PLTS.ABA.GBCA.BindingTrace.value_agree}\leanok
  \uses{def:aba-gbca-spec, thm:gbca-spec-binding, thm:trace-support}
  $\mathrm{BoundTrace}$ says that any two round-$r$ returns of a trace announce the same bit, and $\mathrm{BindingTrace}$ adds that every round-$r$ return of the trace which hands out a value hands out that bit.
  Every trace in the support of every achievable trace distribution of the round-$r$ GBCA specification instance satisfies $\mathrm{BindingTrace}$.\lean{PLTS.ABA.GBCA.specInst_binding}\leanok{} Graded agreement is then a consequence of the trace predicate alone: two returns of the trace that hand out bits each hand out the bit they announce, and the two announcements agree.\lean{PLTS.ABA.GBCA.BindingTrace.value_agree}\leanok{}
\end{theorem}

\begin{proof}\leanok
  \uses{def:aba-gbca-spec, thm:gbca-spec-binding, thm:trace-support}
  Theorem~\ref{thm:trace-support} turns a positive-probability trace into a genuine execution of the instance in which each round-$r$ return label of the trace is a step of a return rule at some state.
  The two clauses are then Theorem~\ref{thm:gbca-spec-binding} read at those steps: the announcements of any two of them agree, and a value-bearing one announces the value it hands out.
  Nothing beyond positive probability is asked of the trace, and nothing beyond membership in one execution of the returns.
  This is what carries binding down a chain of refinements: the property is stated on labels, so an inclusion of achievable trace distributions transports it.
\end{proof}
